\documentclass[aspectratio=169]{beamer}

\usetheme{Madrid}
\usecolortheme{default}

\usepackage[spanish]{babel}
\usepackage[T1]{fontenc}
\usepackage{lmodern}
\usepackage{amsmath,amssymb}
\usepackage{array}
\usepackage{xcolor}

\definecolor{USTGreen}{RGB}{27,94,32}

\setbeamercolor{structure}{fg=USTGreen}
\setbeamercolor{frametitle}{fg=white,bg=USTGreen}
\setbeamercolor{title}{fg=white,bg=USTGreen}
\setbeamertemplate{navigation symbols}{}

\title[Sistemas lineales]{Sistemas de Ecuaciones Lineales con Dos Incognitas}
\subtitle{Clasificacion, metodos algebraicos e interpretacion}
\author{Fundamentos de Matematica Basica}
\institute{Universidad Santo Tomas}
\date{}

\begin{document}
	
	\begin{frame}
		\titlepage
	\end{frame}
	
	\begin{frame}{Objetivo general}
		Analizar y resolver sistemas de ecuaciones lineales con dos incognitas mediante metodos algebraicos y geometricos.
		
		\begin{block}{Resultados de aprendizaje}
			\begin{itemize}
				\item Identificar tipos de sistemas lineales.
				\item Resolver sistemas por sustitucion, igualacion y reduccion.
				\item Interpretar geometricamente las soluciones.
				\item Detectar inconsistencias algebraicas.
			\end{itemize}
		\end{block}
	\end{frame}
	
	\begin{frame}{Definicion}
		Un sistema de ecuaciones lineales con dos incognitas tiene la forma:
		
		\[
		\begin{cases}
			ax+by=c\\
			dx+ey=f
		\end{cases}
		\]
		
		donde \(a,b,c,d,e,f \in \mathbb{R}\).
		
		\bigskip
		
		Resolver el sistema significa encontrar los valores de \(x\) e \(y\) que satisfacen ambas ecuaciones.
	\end{frame}
	
	\begin{frame}{Interpretacion geometrica}
		Cada ecuacion lineal representa una recta en el plano cartesiano.
		
		\bigskip
		
		\begin{center}
			\begin{tabular}{|c|c|c|}
				\hline
				\textbf{Tipo} & \textbf{Rectas} & \textbf{Soluciones} \\
				\hline
				Compatible determinado & Se cortan & Una solucion \\
				\hline
				Compatible indeterminado & Coinciden & Infinitas soluciones \\
				\hline
				Incompatible & Paralelas & Sin solucion \\
				\hline
			\end{tabular}
		\end{center}
	\end{frame}
	
	\begin{frame}{Sistema compatible determinado}
		\[
		\begin{cases}
			x+y=5\\
			2x-y=4
		\end{cases}
		\]
		
		Sumamos:
		
		\[
		3x=9
		\]
		
		\[
		x=3
		\]
		
		Luego:
		
		\[
		3+y=5
		\]
		
		\[
		y=2
		\]
		
		\[
		\boxed{(x,y)=(3,2)}
		\]
	\end{frame}
	
	\begin{frame}{Sistema compatible indeterminado}
		\[
		\begin{cases}
			x+y=4\\
			2x+2y=8
		\end{cases}
		\]
		
		La segunda ecuacion es multiplo de la primera:
		
		\[
		2(x+y)=8
		\]
		
		\[
		x+y=4
		\]
		
		\[
		\boxed{\text{Infinitas soluciones}}
		\]
	\end{frame}
	
	\begin{frame}{Sistema incompatible}
		\[
		\begin{cases}
			x+y=3\\
			x+y=7
		\end{cases}
		\]
		
		La misma expresion no puede tomar dos valores distintos.
		
		\[
		\boxed{\text{No tiene solucion}}
		\]
	\end{frame}
	
	\begin{frame}{Metodo de sustitucion}
		\[
		\begin{cases}
			x+y=10\\
			x=y+2
		\end{cases}
		\]
		
		Sustituimos:
		
		\[
		(y+2)+y=10
		\]
		
		\[
		2y=8
		\]
		
		\[
		y=4
		\]
		
		\[
		x=6
		\]
		
		\[
		\boxed{(x,y)=(6,4)}
		\]
	\end{frame}
	
	\begin{frame}{Metodo de igualacion}
		\[
		\begin{cases}
			x=8-y\\
			x=2y+2
		\end{cases}
		\]
		
		Igualamos:
		
		\[
		8-y=2y+2
		\]
		
		\[
		6=3y
		\]
		
		\[
		y=2
		\]
		
		\[
		x=6
		\]
		
		\[
		\boxed{(x,y)=(6,2)}
		\]
	\end{frame}
	
	\begin{frame}{Metodo de reduccion}
		\[
		\begin{cases}
			2x+y=11\\
			x-y=1
		\end{cases}
		\]
		
		Sumamos:
		
		\[
		3x=12
		\]
		
		\[
		x=4
		\]
		
		\[
		4-y=1
		\]
		
		\[
		y=3
		\]
		
		\[
		\boxed{(x,y)=(4,3)}
		\]
	\end{frame}
	
	\begin{frame}{Ejercicio 1}
		Resolver por sustitucion:
		
		\[
		\begin{cases}
			x+y=18\\
			x=2y
		\end{cases}
		\]
		
		\bigskip
		
		Indicar:
		
		\begin{itemize}
			\item solucion;
			\item tipo de sistema;
			\item verificacion.
		\end{itemize}
	\end{frame}
	
	\begin{frame}{Ejercicio 2}
		Resolver por igualacion:
		
		\[
		\begin{cases}
			x=20-2y\\
			x=y+5
		\end{cases}
		\]
		
		\bigskip
		
		Clasificar el sistema e interpretar geometricamente.
	\end{frame}
	
	\begin{frame}{Ejercicio 3}
		Resolver por reduccion:
		
		\[
		\begin{cases}
			3x+2y=24\\
			x-2y=0
		\end{cases}
		\]
	\end{frame}
	
	\begin{frame}{Patrones e inconsistencias}
		\begin{block}{Contradiccion}
			Si aparece una igualdad falsa:
			
			\[
			0=5
			\]
			
			el sistema es incompatible.
		\end{block}
		
		\begin{block}{Identidad}
			Si aparece una igualdad verdadera:
			
			\[
			0=0
			\]
			
			el sistema es compatible indeterminado.
		\end{block}
	\end{frame}
	
	\begin{frame}{Cierre}
		\begin{itemize}
			\item Un sistema lineal con dos incognitas representa dos rectas.
			\item Puede tener una solucion, infinitas soluciones o ninguna.
			\item Los metodos principales son sustitucion, igualacion y reduccion.
			\item Toda solucion debe verificarse.
		\end{itemize}
	\end{frame}
	
\end{document}